\scanpage{106}
Donc \deleted{\unsure{en posant|remplaçant}} pour $p+q=i-1$, on trouve un poids $\le 2(n-1)-i+1=2n-i-1$.

\uline{Remarque} \quad La filtration de $H^*(Y)$ déduite de la suite spectrale \unsure{$\infty$|SS|ci} n'est autre que la filtration par les poids. La filtration image de $H^i_!$ est donc aussi la filtration par les poids croissants (de $0$ à $i$, le dernier cran \unsure{comme|connu} noyau) donc est également indépendante de la \unsure{compactification} $X'$ choisie. Écrivons les termes de la filtration, en utilisant le fait que $H^{i-1}(X')$ est purement de poids $i-1$, et \deleted{$E_2^{i-1,0}$} que la suite spectrale $\infty$ dégénère pour des raisons de poids :

\diagram{$\begin{array}{c|c|c|c|c|c|c}
\hline
 &E_2^{i-1,0},&E_2^{i-2,1},&\ldots&E_2^{1,i-2},&
 E_2^{0,i-1}/\operatorname{Im}H^{i-1}(X'),&
 \operatorname{Ker}(H^i(X')\to E_2^{0,i})\\ \hline
 \text{poids}&0&1&&i-2&i-1&i\\ \hline\hline
 \deleted{\text{\unsure{niveau}}}&&&&&&
\end{array}$}

\deleted{\unsure{Car si on|Cas $i\ge n$}} Si $i\le n$, les \textquotedblleft{}niveaux sont \inserted{\unsure{seult|seulm.}} $\le$ poids, il n'y a pas lieu à priori d'attendre plus. L'on trouve d'autre part, pour $i\ge n$, un tableau

\diagram{$\begin{array}{c|c|c|c|c|c|c}
 &\overset{\text{\unsure{degré filtrant|de gr. filtré} }\ge n}{0\quad\ldots\quad0}
 &E_2^{\substack{\deleted{\unsure{2n-i,\,i-1}}\\2n-i-1,2(i-n)}}
 &E_2^{2n-i-2,2(i-n)+1}&\ldots
 &E_2^{0,i-1}/\operatorname{Im} H^{i-1}(X')&
 \operatorname{Ker}(H^i(X')\to E_2^{0,\unsure{i}})\\ \hline
 \text{poids}&\smash{\raisebox{-.7ex}{$\times$}}&2(i-n)&2(i-n)+1&\ldots&i-1&i\\ \hline
 \text{niveau}&&0&1&\ldots&2n-i-1&2n-i
\end{array}$}

\uline{Corollaire 1} \quad (Par dualité) $H^i(X)$ est de poids compris entre $i$ et $2n$\deleted{\unsure{$-i$}}. Le morceau \deleted{\unsure{pur}} de poids $i$ est exactement l'image de $H^i(X')$. Si $i\deleted{\unsure{<}}\ge n$, les niveaux des morceaux de poids $i,i+1,\ldots,2n$ sont resp. $\le 2n-i,2n-i-1,\ldots,0$. Si $i\le n$, il n'y a de morceaux \deleted{\unsure{de}} que des poids \deleted{$i,i+1,\ldots,2n-2(n-i)=2n$\unsure{ $2i$, \ldots, dont}} \inserted{$i,i+1,\ldots,2i$} qui sont de niveaux resp. $\le i,i-1,\ldots,0$.

% slanted
\margin{Tous les morceaux de la filtration sont \textquotedblleft\uline{effectifs}\textquotedblright}
% slanted
\margin{En tous cas niveau $\le\operatorname{Inf}(i,n)$}
\margin{\uline{NB} les niveaux ici \uline{décroissent}}
% archive sheet number: 105
